% !TEX root = userguide.tex

\def\headdate{12th July 2021}

\section{Stabilized elements for the Stokes problem}
\label{sec:stabilizedstokes}
Various stabilization techniques exist for velocity/pressure elements that are
not LBB stable (such as equal order elements). Most of these stabilization
techniques are based on the
residual. Drawbacks of that approach are:
\begin{enumerate}
 \item The residual contains second-order derivatives.
 \item If the momentum balance contains other terms, like in Navier-Stokes or
       viscoelastic flows, these terms need to be taken into
       account in the residual.
 \item Tweaking parameters need to be introduced for
       which it is not straightforward what values need to be taken.
       Usually, these depend on the relative magnitude of the various terms and
       the size of the elements.
  \item The system becomes non-symmetric, even for the symmetric Stokes system.
\end{enumerate}
In \cite{Dohrmann04}, Dohrmann \& Bochev introduced a simple stabilization
technique based on local pressure projections (and thus does not use the
residual). Only
a least-squares pressure term is added to the weak form.
The additional term can be
computed using standard finite element techniques, adds only to the zero
diagonal block, does not need a tweaking parameter and is symmetric.

For the Dohrmann-Bochev stabilization we need to introduce an additional
pressure finite element space $P'_h$, consisting of piecewise discontinuous
polynomials:
\[
     p_h'=\tsum_i \psi'_i p'_i
\]
in addition to the normal continuous pressure finite element space $P_h$
\[
     p_h=\tsum_k \psi_k p_k
\]
Now we consider the projection $\Pi p_h\in P'_h$ of $p_h\in P_h$ on $P'_h$:
\[
    (q'_h,\Pi p_h-p_h) \text{\quad for all }q_h'\in P'_h
\]
Writing $\Pi p_h= \psi'_i p'_i$ (omitting the $\tsum$ over $i$), we find
\[
    p'_i=m^{-1}_{ij}(\psi_j',p_h)
\]
and thus
\[
   \Pi p_h = \psi'_i m^{-1}_{ij}(\psi_j',p_h)=
                \psi'_i m^{-1}_{ij}(\psi_j',\psi_k)p_k
\]
where $m_{ij}=(\psi'_i,\psi'_j)$ is the pressure mass matrix of $P'_h$. Note,
that $\psi'_i$ is discontinuous across element boundaries and the inverse
$m^{-1}_{ij}$ matrix can be computed on element level.

The Stokes problem is determined by the stationary value of the functional
(saddle point problem)
\[
   L(\vek u,p) = \frac12\big(\ten D(\vek u),\eta\ten D(\vek u)\big)-
     (\nabla\cdot\vek u,p) -(\vek u,\vek f)
\]
The Dohrmann-Bochev stabilization consists of subtracting a least-squares term
from the functional $L(\vek u,p)$:
\[
  L'(\vek u,p) = \frac12\big(\ten D(\vek u),\eta\ten D(\vek u)\big)-
     (\nabla\cdot\vek u,p) -(\vek u,\vek f)-
     \frac1{2\eta'}(p-\Pi p,p-\Pi p)
\]
where $\eta'$ is a viscosity parameter for which $\eta'=\eta$ seems to work
just fine. The least-squares term leads to an additional term in the weak
 form of
\[
     \dots -\frac1{\eta'}(q_h-\Pi q_h,p_h-\Pi p_h) = 0
\]
which adds to the continuity equation and fills the ``zero diagonal block''.
Note, that 
\[
    (\Pi q_h,p_h-\Pi p_h)=0 \text{\quad or\quad} (\Pi q_h,\Pi p_h)=(\Pi q_h,p_h)
\]
because the difference with respect to the projection should be orthogonal to
the projection itself. A different proof is by substituting the expressions for
$\Pi p_h$ and $m_{ij}$:
\begin{multline*}
  (\Pi q_h,\Pi p_h)=
   (\psi'_i m^{-1}_{ij}(\psi_j',q_h), \psi'_k m^{-1}_{km}(\psi_m',p_h))=\\
   (\psi'_i,\psi'_k) m^{-1}_{ij}(\psi_j',q_h)m^{-1}_{km}(\psi_m',p_h)=
        m_{ik} m^{-1}_{ij}(\psi_j',q_h)m^{-1}_{km}(\psi_m',p_h)=\\
        \delta_{kj}(\psi_j',q_h)m^{-1}_{km}(\psi_m',p_h)=
                   (\psi_k',q_h)m^{-1}_{km}(\psi_m',p_h)=(\Pi q_h,p_h)
\end{multline*}
Finally the additional term becomes:
\[
   -\frac1{\eta'}[(q_h,p_h)-(q_h,\Pi p_h)]
\]
with additional matrix
\[
  -\frac1{\eta'}[(\psi_k,\psi_m)-(\psi_k,\psi'_i)m^{-1}_{ij}(\psi'_j,\psi_m)]
\]
which is \emph{symmetric} and negative definite. The structure of the system
matrix becomes
\[
 \begin{bmatrix}
   \mat{S} &-\mat{L}^T \\[1ex]
   -\mat{L} &-\mat C
 \end{bmatrix}
 \begin{bmatrix}
    \col{u}\\[1ex]
    \col p
 \end{bmatrix}
\]
with $\mat C$ symmetric positive definite.

Remarks:
\begin{enumerate}
   \item For equal order elements the polynomial order for $\psi'$ is chosen
   one order less than $\psi$. 
   Even lower order is possible (like a constant for quadratic elements), but
   then accuracy is lost.
   \item The $P_1/P_1$ pair is an attractive alternative for the mini element
   ($P_1^+/P_1)$, with less velocity degrees of freedom and constant gradients
   inside the element.
   \item In \cite{Dohrmann04} it is mentioned that the stabilization might be
   beneficial for iterative solvers, but this has not been tested.
   \item For viscoelastic flow the use of these elements is questionable, for
   example the $Q_2/Q_2$ pair is unstable. Combining the LBB-stable element
   $Q_2/Q_1$ with $P_1$ for $\psi'$ is stable for viscoelastic flow. Using 
   $P_0$ for $\psi'$ is even more stable, but also less accurate.
   \item The procedure is a way to produce equal-order spectral elements.
   Such elements might be used together with finite-element
   preconditioners using stabilized $Q_1/Q_1$ low-order elements (not tested).
 \item The current implementation is quite general, using Cholesky
   $LL^T$ decomposition routines from Lapack. For special cases, like spectral
   elements, it might be more efficient to use/implement special elements.
\end{enumerate}

The following examples use stabilized elements:\\\medskip
\begin{center}
 \begin{tabular}{ll}
  file& elements \\ \hline
\texttt{stokes35.f90} &$Q_2/Q_2$ \\
\texttt{stokes36.f90} &$P_1/P_1$ \\
\texttt{stokes37.f90} &$Q_1/Q_1$ \\
\texttt{stokes38.f90} &$Q_p/Q_p$ spectral elements\\
\texttt{stokes53.f90} &$P_p/P_p$ high-order triangular elements\\
\texttt{stokes54.f90} &$Q_p/Q_p$ high-order quadrilateral elements\\
\texttt{sphere29.f90} &$Q_p/Q_p$ high-order quadrilateral elements\\
\texttt{sphere30.f90} &$P_p/P_p$ high-order triangular elements
\end{tabular}
\end{center}
 
